\documentclass{MFM-hw} % You must have the MFM5210.cls file in your project or working directory (i.e. the same directory as this document) 

\HWauthor{Wei Qi}{226010267@link.cuhk.edu.cn} % Put your name and CUHK-SZ e on this line
\HWno{1} % Enter the number of the homework you are working on
\HWcourse{MFM 5210} % Enter the course department and number here
\date{Sep 14, 2026}
\usepackage{amsmath}
\usepackage{amsfonts}
\usepackage{bbm}
\usepackage{hyperref}
\newcommand{\E}{\operatorname{E}}
\newcommand{\N}{\mathcal{N}}
\newcommand{\dd}{\,\mathrm{d}}
\newcommand{\EE}{\mathbb{E}}
\newcommand{\PP}{\mathbb{P}}


\begin{document}
\maketitle
\HWproblem
\HWsubproblem
Which of the following events belong to the $\sigma$-algebra $\mathcal{F}_2$?
\[
    \{B_1 \leq 2\},\quad \{|B_2 - B_1| \geq 3\}, \quad \{B_1 \neq B_2\}, \quad\{B_3 \geq 0\}, \quad \{\int_0^3 B_t \mathrm{d}t < 10\}.
\]

$\{B_1 < 2\} \in \mathcal{F}_2$, for $B_1$ is $\mathcal{F}_1$-measurable and $\mathcal{F}_1 \subset \mathcal{F}_2$. ${|B_2 - B_1| < 2} \in \mathcal{F}_2$, for $B_1, B_2$ are both $\mathcal{F}_2$-measurable. $\{B_1 \neq B_2\} \in \mathcal{F}_2$ for the same reason.$\{B_3 \geq 0\} \not \in \mathcal{F}_2$ for it depends on $B_3$, which is independent of $\mathcal{F}_2$. $\{\int_0^3 B_t \mathrm{d}t < 10\} \not \in \mathcal{F}_2$. The integral uses the path on $[0,3]$, including times after 2, so it is $\mathcal{F}_3$-measurable but not $\mathcal{F}_2$-measurable.

To conlude, the following events belong to $\mathcal{F}_2$
\[
\{B_1 \leq 2\},\quad \{|B_2 - B_1| \geq 3\}, \quad \{B_1 \neq B_2\}.
\]
while others do not.

\HWsubproblem
Which of the following random variables are $\mathcal{F}_3$-measurable?
\[
    (B_1 + B_2)^2, \quad \sin B_3, \quad B_4 - B_3, \quad \mathbb{E}[B_5^2], \quad \max\limits_{1\leq t\leq 3}B_t, \quad \int_0^3 e^{B_t} \mathrm{d}t.
\]

Because
\[
\mathcal{F}_3 = \sigma\{X_t : t \leq 3\} 
\]
contains information available from past to time 3 and $B_4, B_5$ are independent of $\mathcal{F}_3$, the following random variables are  $\mathcal{F}_3$-measurable
\[
    (B_1 + B_2)^2, \quad \sin B_3, \quad \max\limits_{1\leq t\leq 3}B_t, \quad \int_0^3 e^{B_t} \mathrm{d}t.
\]
while others are not.

\HWproblem
Fix $0<s<t$.

\HWsubproblem
What is the distribution of $B_s + B_t$?

\textbf{Lemma:} Let $X$ and $Y$ be independent random variables that are normally distributed. If
\[
X \sim N(\mu_X, \sigma^2_X), \quad Y \sim N(\mu_Y, \sigma^2_Y),
\]
their sum $Z = X + Y$ satisfies
\[
Z \sim N(\mu_X+\mu_Y, \sigma^2_X+\sigma^2_Y).
\]

\textbf{Proof of lemma:} 

For independent random variables $X$ and $Y$, the distribution $f_Z$ of $Z = X + Y$ equals the
convolution of $f_X$ and $f_Y$:
\begin{equation*}
f_Z(z) = \int_{-\infty}^{\infty} f_Y(z - x)\, f_X(x) \dd x .
\end{equation*}

Given that $f_X$ and $f_Y$ are normal densities,
\begin{equation*}
\begin{aligned}
f_X(x) &= \N(x; \mu_X, \sigma_X^2) = \frac{1}{\sqrt{2\pi}\,\sigma_X} e^{-(x-\mu_X)^2/(2\sigma_X^2)} \\[5pt]
f_Y(y) &= \N(y; \mu_Y, \sigma_Y^2) = \frac{1}{\sqrt{2\pi}\,\sigma_Y} e^{-(y-\mu_Y)^2/(2\sigma_Y^2)}
\end{aligned}
\end{equation*}

Substituting into the convolution:
\begin{align*}
f_Z(z)
&= \int_{-\infty}^{\infty}
   \frac{1}{\sqrt{2\pi}\,\sigma_Y}
   \exp \left[ -\frac{(z-x-\mu_Y)^2}{2\sigma_Y^2} \right]
   \frac{1}{\sqrt{2\pi}\,\sigma_X}
   \exp \left[ -\frac{(x-\mu_X)^2}{2\sigma_X^2} \right]
   \dd x \\[6pt]
&= \int_{-\infty}^{\infty}
   \frac{1}{\sqrt{2\pi}\sqrt{2\pi}\,\sigma_X \sigma_Y}
   \exp \left[ -\frac{\sigma_X^2 (z-x-\mu_Y)^2 + \sigma_Y^2 (x-\mu_X)^2}{2\sigma_X^2 \sigma_Y^2} \right]
   \dd x \\[6pt]
&= \int_{-\infty}^{\infty}
   \frac{1}{\sqrt{2\pi}\sqrt{2\pi}\,\sigma_X \sigma_Y}
   \exp \left[ -\frac{
      \begin{aligned}
        &\sigma_X^2 \left(z^2 + x^2 + \mu_Y^2 - 2xz - 2z\mu_Y + 2x\mu_Y\right) \\
        &\quad {}+ \sigma_Y^2 \left(x^2 + \mu_X^2 - 2x\mu_X\right)
      \end{aligned}
   }{2\sigma_Y^2 \sigma_X^2} \right]
   \dd x \\[6pt]
&= \int_{-\infty}^{\infty}
   \frac{1}{\sqrt{2\pi}\sqrt{2\pi}\,\sigma_X \sigma_Y}
   \exp \left[ -\frac{
      \begin{aligned}
        &x^2 \left(\sigma_X^2 + \sigma_Y^2\right) -
         2x\left(\sigma_X^2 (z - \mu_Y) + \sigma_Y^2 \mu_X\right) \\
        &\quad {}+ \sigma_X^2 \left(z^2 + \mu_Y^2 - 2z\mu_Y\right) + \sigma_Y^2 \mu_X^2
      \end{aligned}
   }{2\sigma_Y^2 \sigma_X^2} \right]
   \dd x
\end{align*}

Defining $\sigma_Z = \sqrt{\sigma_X^2 + \sigma_Y^2}$, and completing the square:
\begin{align*}
f_Z(z)
&= \int_{-\infty}^{\infty}
   \frac{1}{\sqrt{2\pi}\,\sigma_Z}
   \frac{1}{\sqrt{2\pi}\,\frac{\sigma_X \sigma_Y}{\sigma_Z}}
   \exp \left[ -\frac{
      \begin{aligned}
        &x^2 - 2x\,\frac{\sigma_X^2 (z - \mu_Y) + \sigma_Y^2 \mu_X}{\sigma_Z^2} \\
        &\quad {}+ \frac{\sigma_X^2 \left(z^2 + \mu_Y^2 - 2z\mu_Y\right) + \sigma_Y^2 \mu_X^2}{\sigma_Z^2}
      \end{aligned}
   }{2\left(\frac{\sigma_X \sigma_Y}{\sigma_Z}\right)^{2}} \right]
   \dd x \\[6pt]
&= \int_{-\infty}^{\infty}
   \frac{1}{\sqrt{2\pi}\,\sigma_Z}
   \frac{1}{\sqrt{2\pi}\,\frac{\sigma_X \sigma_Y}{\sigma_Z}}
   \exp \left[ -\frac{
      \begin{aligned}
        &\left(x - \frac{\sigma_X^2 (z - \mu_Y) + \sigma_Y^2 \mu_X}{\sigma_Z^2}\right)^{2} \\
        &\quad {} - \left(\frac{\sigma_X^2 (z - \mu_Y) + \sigma_Y^2 \mu_X}{\sigma_Z^2}\right)^{2} \\
        &\quad {} + \frac{\sigma_X^2 (z - \mu_Y)^2 + \sigma_Y^2 \mu_X^2}{\sigma_Z^2}
      \end{aligned}
   }{2\left(\frac{\sigma_X \sigma_Y}{\sigma_Z}\right)^{2}} \right]
   \dd x \\[6pt]
&= \int_{-\infty}^{\infty}
   \frac{1}{\sqrt{2\pi}\,\sigma_Z}
   \begin{aligned}[t]
     &\exp \left[ -\frac{
        \begin{aligned}
          &\sigma_Z^2 \left(\sigma_X^2 (z - \mu_Y)^2 + \sigma_Y^2 \mu_X^2\right) \\
          &\quad {} - \left(\sigma_X^2 (z - \mu_Y) + \sigma_Y^2 \mu_X\right)^{2}
        \end{aligned}
     }{2\sigma_Z^2 (\sigma_X \sigma_Y)^{2}} \right] \\[3pt]
     &\quad \times \frac{1}{\sqrt{2\pi}\,\frac{\sigma_X \sigma_Y}{\sigma_Z}}
       \exp \left[ -\frac{
          \left(x - \frac{\sigma_X^2 (z - \mu_Y) + \sigma_Y^2 \mu_X}{\sigma_Z^2}\right)^{2}
       }{2\left(\frac{\sigma_X \sigma_Y}{\sigma_Z}\right)^{2}} \right]
   \end{aligned}
   \dd x \\[6pt]
&= \frac{1}{\sqrt{2\pi}\,\sigma_Z}
   \exp \left[ -\frac{(z - (\mu_X + \mu_Y))^2}{2\sigma_Z^2} \right] \notag \\[3pt]
&\quad \times \int_{-\infty}^{\infty}
   \frac{1}{\sqrt{2\pi}\,\frac{\sigma_X \sigma_Y}{\sigma_Z}}
   \exp \left[ -\frac{
      \left(x - \frac{\sigma_X^2 (z - \mu_Y) + \sigma_Y^2 \mu_X}{\sigma_Z^2}\right)^{2}
   }{2\left(\frac{\sigma_X \sigma_Y}{\sigma_Z}\right)^{2}} \right]
   \dd x
\end{align*}

The expression in the integral is a normal density distribution on $x$, and therefore the integral
is evaluated to $1$. The desired result follows
\begin{equation*}
f_Z(z) = \frac{1}{\sqrt{2\pi}\,\sigma_Z}
         \exp \left[ -\frac{(z - (\mu_X + \mu_Y))^2}{2\sigma_Z^2} \right]
\end{equation*}
which proves that
\[
Z \sim N(\mu_X+\mu_Y, \sigma^2_X+\sigma^2_Y).
\]

\noindent
(Source of the proof:
\url{https://en.wikipedia.org/wiki/Sum_of_normally_distributed_random_variables}.)

According to the properties of standard Brownian motion, 
\[
B_s \sim N(0,s), \quad (B_t - B_s) \sim N (0, t-s),
\]
and $B_s, B_t - B_s$ are independent, which derive that $2B_s \sim N(0, 4s)$ and $2B_s, B_t - B_s$ are independent. According to the lemma,
\[
B_s + B_t = 2B_s + (B_s - B_t) \sim N(0, 3s + t),
\]
which means $B_s + B_t$ satisfies normal distribution with mean $0$ and variance $3s + t$.

\HWsubproblem
Compute $\PP\{B_s > 0, B_t - B_s > 0\}$.

$B_s$ and $B_t - B_s$ are independent and satisfy
\[
B_s \sim N(0,s), \quad (B_t - B_s) \sim N (0, t-s).
\]
Hence
\[
\PP\{B_s > 0, B_t - B_s > 0\} = \PP\{B_s > 0\}\PP\{B_t - B_s > 0\} = \frac12 \times \frac12 = \frac14.
\]

\HWsubproblem
Compute $\EE[B_2^3B_4]$.

\textbf{Lemma:} Let $\{B_t\}_{t \geq 0}$ be a standard Brownian motion, the $n^{\mathrm{th}}(n \geq 1, n \text{ is integer})$ moment of $B_t$ satisfies
\[
    \EE[B_t] = 
    \begin{cases}
    0, & n \text{ is odd,} \\
    (n-1)!t^{n/2}, & n \text{ is even}.
    \end{cases}
\]

The lemma is proved in class. And $B_2$ and $B_4 - B_2$ are independent, satisfying
\[
B_2 \sim N(0,2), \quad (B_4 - B_2) \sim N (0, 2).
\]
So
\begin{align*}
    \EE[B_2^3B_4] &= \EE[B_2^3(B_4-B_2+B_2)] = \EE[B_2^3(B_4-B_2) + B_2^4] \\
                  &= \EE[B_2^3] \EE[B_4-B_2] +\EE[B_2^4],
\end{align*}
where according to the lemma
\[
    \EE[B_2^3]=0, \quad \EE[B_2^4] = 12,
\]
hence
\[
    \EE[B_2^3B_4] = \EE[B_2^4] = 12.
\]

\HWsubproblem
Compute $\EE[B_1^4B_3^2]$.

$B_1$ and $B_3 - B_1$ are independent, satisfying
\[
B_1 \sim N(0,1), \quad (B_3 - B_1) \sim N (0, 2).
\]
So
\begin{align*}
    \EE[B_1^4B_3^2] &= \EE[B_1^4(B_3-B_1+B_1)^2] = \EE[B_1^4(B_3-B_1)^2+2B_1^5(B_3-B_1)+B_1^6] \\
                    &= \EE[B_1^4]\EE[(B_3-B_1)^2] + 2\EE[B_1^5]\EE[B_3-B_1] + \EE[B_1^6] \\
\end{align*}
where
\[
    \EE[B_1^4] = 3, \quad \EE[(B_3-B_1)^2] = 2, \quad \EE[B_1^5] = \EE[B_3-B-1] = 0, \quad \EE[B_1^6] = 15,
\]
hence
\[
    \EE[B_1^4B_3^2] = 3 \times 2 + 0 + 15 = 21.
\]

\HWsubproblem
Compute $\EE[e^{a(B_t-B_s)}]$ for any $a \in \mathbb{R}$.

Let $X = B_t -B_s \sim N(0, t-s)$, the probability density function of $X$ is 
\[
    f_X(x) = \frac{1}{\sqrt{2\pi (t-s)}} e^{-\frac{x^2}{2(t-s)}},
\]
hence
\begin{align*}
    \EE[e^{a(B_t-B_s)}] &= 
    \EE[e^{a(B_t-B_s)}] = \EE[e^{aX}] = \int_{-\infty}^{+\infty} e^{ax} f_X(x) \dd x  \\
    &= \int_{-\infty}^{+\infty} \frac{1}{\sqrt{2\pi (t-s)}} e^{-\frac{x^2}{2(t-s)}+ax} \dd x \\
    &= \int_{-\infty}^{+\infty} \frac{1}{\sqrt{2\pi (t-s)}} e^{-\frac{[x-a(t-s)]^2}{2(t-s)}+\frac{(t-s)a^2}{2}} \dd x \\
    &= e^{\frac{(t-s)a^2}{2}}\int_{-\infty}^{+\infty} \frac{1}{\sqrt{2\pi (t-s)}} e^{-\frac{[x-a(t-s)]^2}{2(t-s)}} \dd x
\end{align*}

The expression in the integral is a normal density distribution, and therefore the integral
is evaluated to 1. The desired result follows
\[
    \EE[e^{a(B_t-B_s)}] = e^{\frac{(t-s)a^2}{2}}.
\]

\HWproblem
Let $\phi(t)_{t\in[0,5]}$ be the simple process defined by
\[
    \phi_t = \mathbbm{1}_{0<t\leq 1} + |B_1|^3 \mathbbm{1}_{1 < t \leq 3} - B_2 e^{-B_3} \mathbbm{1}_{3<t \leq 5}.
\]

\HWsubproblem
Is $\phi_t$ predictable? Explain your answer.

Since
\[
    a_1 = \mathbbm{1} \in \mathcal{F}_0, \quad a_2 = |B_1|^3 \in \mathcal{F}_1, \quad a_3 = -B_2 e^{-B_3} \in \mathcal{F}_3,
\]
$\phi_t$ is predictable.

\HWsubproblem
Write down the definition of the Itô integral $\int_0^5 \phi_t \dd B_t$. 

The Itô integral $\int_0^5 \phi_t \dd B_t$ is defined as
\[
    \int_0^5 \phi_t \dd B_t = (B_1 - B_0) + |B_1|^3 (B_3 - B_1) - B_2 e^{-B_3}(B_5-B_3).
\]

\HWproblem
Fix $t>0$. Compute the mean and variance of each of the following integrals.

\HWsubproblem
$\int_0^{T} B_t^2 \dd B_t$.

First verify that $\{B_{t}^2\}_{t\in[0,t]} \in V(0,t)$. Because $B_t \in \mathcal{F}_t$ for any $t \in [0,t]$, $B_t^2 \in \mathcal{F}_t$ for any $t \in [0,t]$. Additionally, 
\[
    \EE[\int_0^T |B_t^2|^2 \dd t] = \int_0^T \EE[B_t^4]\dd t = \int_0^{T} 3t^2 \dd t = T^3 < +\infty,
\]
so $\{B_{t}^2\}_{t\in[0,t]} \in V(0,t)$. According to properties of Itô integral,
\[
    \EE[\int_0^{T} B_t^2 \dd B_t] = 0, \quad Var[\int_0^{T} B_t^2 \dd B_t] = \EE[(\int_0^{T} B_t^2 \dd B_t)^2] = \EE[\int_0^T |B_t^2|^2 \dd t] = T^3.
\]

\HWsubproblem
$\int_0^{T} 2e^{B_t} \dd B_t$

First verify that $\{2e^{B_t}\}_{t\in[0,t]} \in V(0,t)$. Because $B_t \in \mathcal{F}_t$ for any $t \in [0,t]$, $2e^{B_t} \in \mathcal{F}_t$ for any $t \in [0,t]$. Additionally, 
\[
    \EE[\int_0^T |2e^{B_t}|^2 \dd t] = \int_0^T \EE[4e^{2B_t}]\dd t = 4\int_0^{T} \EE[e^{2B_t}]\dd t
\]
where according to the moment generating function
\[
    \EE[e^{2B_t}] = e^{2t},
\]
hence
\[
    \EE[\int_0^T |2e^{B_t}|^2 \dd t] = 4\int_0^{T} e^{2t} \dd t = 2e^{2T} - 2 < +\infty,
\]
so $\{B_{t}^2\}_{t\in[0,t]} \in V(0,t)$. According to properties of Itô integral,
\[
    \EE[\int_0^{T} 2e^{B_t} \dd B_t] = 0, \quad Var[\int_0^{T}2e^{B_t} \dd B_t] = \EE[(\int_0^{T} 2e^{B_t} \dd B_t)^2] = \EE[\int_0^T |2e^{B_t}|^2 \dd t] = 2e^{2T} - 2.
\]

\HWproblem
Verify that the stochastic process $\{(B_t+t)^2\}_{t \in [0,t]}$ belongs to $V(0,t)$ and compute the mean and variance of $\int_0^{T}(B_t + t)^2 \dd B_t$.

Because $B_t \in \mathcal{F}_t$ for any $t \in [0,t]$, $(B_t+t)^2 \in \mathcal{F}_t$ for any $t \in [0,t]$. Additionally,
\begin{align*}
    \EE[\int_0^T |(B_t+t)^2|^2 \dd t] &= \int_0^T \EE[(B_t+t)^4]\dd t = \int_0^T \EE[B_t^4+4B_t^3t + 6B_t^2t^2 + 4B_t t^3 + t^4]\dd t \\
    & = \int_0^T \EE[B_t^4]\dd t + 4\int_0^T t\EE[B_t^3]\dd t + 6\int_0^T t^2\EE[B_t^2]\dd t + 4\int_0^T t^3 \EE[B_t]\dd t + \int_0^T t^4\dd t\\
    & = \int_0^{T} 3t^2\dd t + 0 + 6\int_0^{T} t^3 \dd t + 0 + \int_0^{T} t^4 \dd t \\
    & = T^3 + \frac32 T^4 + \frac15 T^5 < +\infty,
\end{align*}
so  $\{(B_t+t)^2\}_{t \in [0,t]}$. According to properties of Itô integral,
\[
    \EE[\int_0^{T} (B_t+t)^2 \dd B_t] = 0,
\]
\[
Var[\int_0^{T} (B_t+t)^2 \dd B_t] = \EE[(\int_0^{T} (B_t+t)^2 \dd B_t)^2] = \EE[\int_0^T |(B_t+t)^2|^2 \dd t] = T^3 + \frac32 T^4 + \frac15 T^5.
\]


\end{document}
